\documentclass[12pt]{article}

\input{~/School/Notes/header.tex}

\usepackage{booktabs}

\setlength{\parindent}{1.25cm}
\setlength{\parskip}{0.25em}
\renewcommand{\arraystretch}{1.15}

\title{Лабораторная работа по физике}
\author{Денисов Илья}
\date{}

\begin{document}

    \maketitle

    \vspace{5cm}
    {\centering \LARGE Пружина на весах \par}

    \newpage

    \tableofcontents
    \setcounter{tocdepth}{3}

    \newpage

    \section{Цель работы}

    Измерить зависимость веса пружины, лежащей на весах, от вертикальной координаты её верхнего витка, проверить теоретическую зависимость, определить коэффициент жёсткости одного витка пружины и оценить момент силы при скручивании пружины.

    \section{Задачи}

    \begin{enumerate}
        \item Измерить показания весов при разных положениях верхнего витка пружины.
        \item Пересчитать результаты в $\Delta P = P - P_0$ и $\Delta x = x - x_0$.
        \item Получить теоретическую зависимость между $\Delta P$ и $\Delta x$.
        \item Построить линеаризованный график и определить коэффициент жёсткости одного витка.
        \item Предложить способ оценки момента силы при повороте одного конца пружины на один оборот.
    \end{enumerate}

    \section{Оборудование и метод измерения}

    Использовались пружина <<слинки>>, электронные весы, две деревянные линейки, два деревянных брусочка и малярный скотч. Для уменьшения колебаний показаний весов линейка и пружина фиксировались с помощью брусочков и скотча.

    Координата верхнего витка в исходном состоянии $x_0 = 5$ см. Масса всей пружины по исходным данным $m_0 = 15.45$ г, количество витков $n = 37$. Ускорение свободного падения $g = 9.8$ м/с$^2$.

    Для каждой заданной высоты выполнялись три повторных измерения показаний весов; в таблицу внесены только средние арифметические значения. Индивидуальные показания повторов не приводятся, поскольку для обработки используются средние значения.

    \section{Рабочие формулы}

    Изменение координаты верхнего витка:
    \[
        \Delta x = x - x_0.
    \]

    Изменение показаний весов в единицах массы:
    \[
        \Delta m = m - m_0.
    \]

    Изменение веса пружины:
    \[
        \Delta P = \Delta m \cdot 10^{-3} \cdot g.
    \]

    В ходе опыта верхние витки поднимаются, поэтому показания весов уменьшаются. Далее для проверки теоретической зависимости используем модуль изменения веса $|\Delta P|$.

    Пусть $m_{\text{вит}}$ — масса одного витка, а $k$ — жёсткость одного витка. Если изменили положение $n'$ витков, то изменение показаний весов по модулю равно
    \[
        |\Delta P| = m_{\text{вит}} g n'.
    \]

    Для большого числа изменивших положение витков суммарное изменение высоты приближённо равно
    \[
        \Delta x = \frac{m_{\text{вит}} g (n')^2}{2k}.
    \]

    Исключая $n'$, получаем
    \[
        (\Delta P)^2 = 2 k m_{\text{вит}} g \Delta x.
    \]

    Следовательно, график $(\Delta P)^2(\Delta x)$ должен быть линейным, а его угловой коэффициент $\beta$ равен
    \[
        \beta = 2 k m_{\text{вит}} g,
        \qquad
        k = \frac{\beta}{2 m_{\text{вит}} g}.
    \]

    Масса одного витка:
    \[
        m_{\text{вит}} = \frac{m_0}{n}.
    \]

    \section{Результаты измерений}

    \small
    \begin{table}[H]
        \centering
        \begin{xltabular}{\textwidth}{@{}rrrrrr@{}}
            \toprule
            $N$ & $x$, см & $m$, г & $\Delta x$, см & $\Delta m$, г & $|\Delta P|$, Н\\
            \midrule
            1  & 6  & 11.40 & 1  & -4.05  & 0.03969 \\
            2  & 7  & 10.30 & 2  & -5.15  & 0.05047 \\
            3  & 8  & 9.00  & 3  & -6.45  & 0.06321 \\
            4  & 9  & 8.07  & 4  & -7.38  & 0.07232 \\
            5  & 10 & 7.36  & 5  & -8.09  & 0.07928 \\
            6  & 11 & 6.57  & 6  & -8.88  & 0.08702 \\
            7  & 12 & 5.85  & 7  & -9.60  & 0.09408 \\
            8  & 13 & 5.23  & 8  & -10.22 & 0.10016 \\
            9  & 14 & 4.60  & 9  & -10.85 & 0.10633 \\
            10 & 15 & 3.80  & 10 & -11.65 & 0.11417 \\
            11 & 16 & 3.24  & 11 & -12.21 & 0.11966 \\
            12 & 17 & 2.76  & 12 & -12.69 & 0.12436 \\
            13 & 18 & 2.21  & 13 & -13.24 & 0.12975 \\
            14 & 19 & 1.68  & 14 & -13.77 & 0.13495 \\
            15 & 20 & 1.32  & 15 & -14.13 & 0.13847 \\
            16 & 21 & 0.71  & 16 & -14.74 & 0.14445 \\
            17 & 22 & 0.20  & 17 & -15.25 & 0.14945 \\
            \bottomrule
        \end{xltabular}
        \caption{Средние показания весов и пересчитанные величины}
        \label{tab:measurements}
    \end{table}
    \normalsize

    В таблице $m_0 = 15.45$ г и $x_0 = 5$ см. Величины $\Delta x$ в формулах далее переводятся в метры.

    \section{Проверка теоретической зависимости}

    Для проверки теории строим график $(\Delta P)^2$ от $\Delta x$. По линейной аппроксимации экспериментальных данных получено
    \[
        (\Delta P)^2 = \beta \Delta x + b,
    \]
    где
    \[
        \beta \approx 0.1299 \ \text{Н}^2/\text{м},
        \qquad
        b \approx -6.8 \cdot 10^{-5} \ \text{Н}^2.
    \]

    Свободный член мал по сравнению с максимальными значениями $(\Delta P)^2$, а точки хорошо ложатся на прямую. Коэффициент детерминации составляет примерно $R^2 = 0.999$. Следовательно, в пределах точности измерений эксперимент согласуется с теоретической зависимостью $(\Delta P)^2 \sim \Delta x$.

    Для графика приняты оценки погрешностей $\Delta x=0.005~\text{м}$ и $\Delta m=0.01~\text{г}$ для показаний весов. Для $Y=(\Delta P)^2$ погрешность оценена распространением погрешности:
    \[
        \Delta Y \simeq 2|\Delta P|\,\Delta P,\qquad
        \Delta P=0.01\cdot10^{-3}g.
    \]

    \begin{figure}[H]
        \centering
        \begin{tikzpicture}
            \begin{axis}[
                width=0.9\textwidth,
                height=8cm,
                xlabel={$\Delta x$, м},
                ylabel={$(\Delta P)^2$, Н$^2$},
                xmin=0, xmax=0.18,
                ymin=0, ymax=0.025,
                grid=major,
                ]
                \addplot+[
                    only marks, mark=+, mark size=2.5pt,
                    error bars/x dir=both, error bars/x explicit,
                    error bars/y dir=both, error bars/y explicit
                ] coordinates {
                    (0.01,0.0015752) +- (0.005,0.0000078)
                    (0.02,0.0025472) +- (0.005,0.0000099)
                    (0.03,0.0039955) +- (0.005,0.0000124)
                    (0.04,0.0052304) +- (0.005,0.0000141)
                    (0.05,0.0062856) +- (0.005,0.0000155)
                    (0.06,0.0075725) +- (0.005,0.0000171)
                    (0.07,0.0088510) +- (0.005,0.0000184)
                    (0.08,0.0100312) +- (0.005,0.0000196)
                    (0.09,0.0113061) +- (0.005,0.0000209)
                    (0.10,0.0130348) +- (0.005,0.0000228)
                    (0.11,0.0143184) +- (0.005,0.0000239)
                    (0.12,0.0154660) +- (0.005,0.0000249)
                    (0.13,0.0168355) +- (0.005,0.0000260)
                    (0.14,0.0182104) +- (0.005,0.0000270)
                    (0.15,0.0191769) +- (0.005,0.0000279)
                    (0.16,0.0208664) +- (0.005,0.0000289)
                    (0.17,0.0223353) +- (0.005,0.0000299)
                };
                \addlegendentry{Эксперимент с крестами погрешностей}
                \addplot[very thick, domain=0:0.18] {0.1299278*x - 0.00006777};
                \addlegendentry{Линейная аппроксимация}
            \end{axis}
        \end{tikzpicture}
        \caption{Зависимость $(\Delta P)^2(\Delta x)$ с крестами погрешностей}
        \label{fig:linear}
    \end{figure}

    \section{Коэффициент жёсткости одного витка}

    Масса одного витка:
    \[
        m_{\text{вит}} = \frac{15.45 \cdot 10^{-3}}{37}
        \approx 4.18 \cdot 10^{-4} \ \text{кг}.
    \]

    Из углового коэффициента графика:
    \[
        k = \frac{\beta}{2 m_{\text{вит}} g}
        = \frac{0.1299}{2 \cdot 4.18 \cdot 10^{-4} \cdot 9.8}
        \approx 15.9 \ \text{Н/м}.
    \]

    Для оценки погрешности используем стандартную ошибку углового коэффициента линейной регрессии. По остаткам аппроксимации получаем $\Delta\beta \approx 0.00102~\text{Н}^2/\text{м}$, то есть относительная погрешность наклона составляет около $0.8\%$. Масса пружины и число витков также измерены с конечной точностью; для массы одного витка примем относительную погрешность $3\%$ как консервативную оценку. Тогда для частного в первом приближении:
    \[
        \varepsilon_k \approx \frac{\Delta\beta}{\beta}+
        \frac{\Delta m_{\text{вит}}}{m_{\text{вит}}}
        \approx 0.8\%+3\%=3.8\%.
    \]
    Следовательно, $\Delta k \approx 0.038\cdot15.9\approx0.6~\text{Н/м}$, и
    \[
        k=(15.9\pm0.6)~\text{Н/м}.
    \]
    Эта оценка предполагает, что систематические ошибки измерения координаты и показаний весов не превосходят принятых оценок.

    \section{Оценка момента силы при скручивании}

    Для оценки момента силы можно закрепить один конец пружины на длинной линейке, расположенной горизонтально, и опереть свободный конец линейки на весы. Брусочки по сторонам линейки служат направляющими и уменьшают её боковое смещение. После обнуления весов второй конец пружины поворачивают относительно первого на один оборот, не допуская смещения установки.

    Изменение показания весов $\Delta m_{\text{скр}}$ создаёт силу реакции $F = \Delta m_{\text{скр}} g$. Если расстояние от оси поворота до точки опоры линейки равно $L$, то оценка момента:
    \[
        M = \Delta m_{\text{скр}} \cdot 10^{-3} \cdot g \cdot L.
    \]

    По имеющейся записи изменение показания весов при скручивании составляет $1.5$ г. При оценке примем плечо силы равным длине длинной линейки $L \approx 0.40$ м (если фактически измеренное расстояние до точки опоры отличается, следует подставить его), тогда
    \[
        M = 1.5 \cdot 10^{-3} \cdot 9.8 \cdot 0.40
        \approx 5.9 \cdot 10^{-3} \ \text{Н} \cdot \text{м}.
    \]

    Это оценка порядка величины; её точность ограничена тем, что длина плеча и показание весов известны с ограниченной точностью.

    \section{Вывод}

    \begin{enumerate}
        \item Измерена зависимость показаний весов от вертикальной координаты верхнего витка пружины. Получено $17$ точек, а диапазон изменения координаты составил от $1$ до $17$ см.
        \item После пересчёта данных построен график $(\Delta P)^2(\Delta x)$. Он близок к прямой, что согласуется с теоретической зависимостью $(\Delta P)^2 = 2 k m_{\text{вит}} g \Delta x$.
        \item По угловому коэффициенту графика найдена жёсткость одного витка: $k = (15.9 \pm 1.1)$ Н/м. Относительная оценка погрешности составляет около $7\%$.
        \item Предложен способ оценки момента силы при скручивании пружины. При показании весов $1.5$ г и плече $0.40$ м получается оценка $M \approx 5.9 \cdot 10^{-3}$ Н$\cdot$м.
    \end{enumerate}

\end{document}
